% Ad Atticum 5.20 (ad-atticum-05-20)
% Source: https://raw.githubusercontent.com/PerseusDL/canonical-latinLit/master/data/phi0474/phi057/phi0474.phi057.perseus-lat1.xml
% Fetched by scripts/fetch_latin.py

% Dateline (Perseus): Scr. in Cilicia inter a. d. xii et iv K. Ian. a. 703 (51).

\ciceroLetterOpener{CICERO ATTICO salutem}

\ciceroSection{1} Saturnalibus mane se mihi Pindenissitae dediderunt septimo et quinquagesimo die postquam oppugnare eos coepimus. qui malum! isti Pindenissitae qui sunt? inquies; nomen audivi numquam. quid ego faciam? num potui Ciliciam Aetoliam aut Macedoniam reddere? hoc iam sic habeto nec hoc exercitu nec hic tanta negotia geri potuisse. quae cognosce ἐν ἐπιτομῇ ; sic enim concedis mihi proximis litteris. Ephesum ut venerim nosti, qui etiam mihi gratulatus es illius diei celebritatem qua nihil me umquam delectavit magis. inde in oppidis iis †quae erant† mirabiliter accepti Laodiceam pridie Kal. Sextilis venimus. ibi morati biduum perinlustres fuimus honorificisque verbis omnis iniurias revellimus superiores, quod idem Colossis , dein Apameae quinque dies morati et Synnadis triduum, Philomeli quinque dies, Iconi decem fecimus. nihil ea iuris dictione aequabilius, nihil lenius, nihil gravius.

\ciceroSection{2} inde in castra veni a. d. VII Kalendas Septembris. A. d. III exercitum lustravi apud Iconium. ex his castris, cum graves de Parthis nuntii venirent, perrexi in Ciliciam per Cappadociae partem eam quae Ciliciam attingit, eo consilio ut Armenius Artavasdes et ipsi Parthi Cappadocia se excludi putarent. cum dies quinque ad Cybistra Cappadociae castra habuissem, certior sum factus Parthos ab illo aditu Cappadociae longe abesse, Ciliciae magis imminere. itaque confestim iter in Ciliciam feci per Tauri pylas. Tarsum veni a. d. III Nonas Octobris.

\ciceroSection{3} inde ad Amanum contendi qui Syriam a Cilicia in aquarum divertio dividit; qui mons erat hostium plenus sempiternorum. hic a. d. iii Idus Octobr. magnum numerum hostium occidimus. castella munitissima nocturno Pomptini adventu, nostro matutino cepimus, incendimus. imperatores appellati sumus. castra paucos dies habuimus ea ipsa quae contra Darium habuerat apud Issum Alexander, imperator haud paulo melior quam aut tu aut ego. ibi dies quinque morati direpto et vastato Amano inde discessimus. interim (scis enim dici quaedam πανικά , dici item τὰ κενὰ τοῦ πολέμου ) rumore adventus nostri et Cassio qui Antiochia tenebatur animus accessit et Parthis timor iniectus est. itaque eos cedentis ab oppido Cassius insecutus rem bene gessit. qua in fuga magna auctoritate Osaces dux Parthorum vulnus accepit eoque interiit paucis post diebus. erat in Syria nostrum nomen in gratia.

\ciceroSection{4} venit interim Bibulus; credo, voluit appellatione hac inani nobis esse par. in eodem Amano coepit loreolam in mustaceo quaerere. at ille cohortem primam totam perdidit centurionemque primi pili nobilem sui generis Asinium Dentonem et reliquos cohortis eiusdem et Sex. Lucilium, T. Gavi Caepionis locupletis et splendidi hominis filium, tribunum militum. sane plagam odiosam acceperat cum re tum tempore. nos ad Pindenissum, quod oppidum munitissimum Eleutherocilicum omnium memoria in armis fuit. feri homines et acres et omnibus rebus ad defendendum parati. cinximus vallo et fossa; aggere maximo, vineis, turre altissima, magna tormentorum copia, multis sagittariis, magno labore, apparatu multis sauciis nostris, incolumi exercitu negotium confecimus. hilara sane Saturnalia militibus quoque quibus equis exceptis reliquam praedam concessimus. mancipia venibant Saturnalibus tertiis. cum haec scribebam, in tribunali res erat ad HS c_x_x_ . hinc exercitum in hiberna agri male pacati deducendum Quinto fratri dabam; ipse me Laodiceam recipiebam.

\ciceroSection{6} haec adhuc. sed ad praeterita revertamur. quod me maxime hortaris et quod pluris est quam omnia, in quo laboras ut etiam Ligurino Μώμῳ satis faciamus, moriar si quicquam fieri potest elegantius. nec iam ego hanc continentiam appello, quae virtus voluptati resistere videtur. ego in vita mea nulla umquam voluptate tanta sum adfectus quanta adficior hac integritate, nec me tam fama quae summa est quam res ipsa delectat. quid quaeris? fuit tanti. me ipse non noram nec satis sciebam quid in hoc genere facere possem. recte πεφυσίωμαι . nihil est praeequis carius . interim haec λαμπρά . Ariobarzanes opera mea vivit, regnat; ἐν consilio et auctoritate et quod insidiatoribus eius ἀπρόσιτον me non modo ἀδωροδόκητον praebui regem regnumque servavi. interea e Cappadocia ne pilum quidem. Brutum abiectum quantum potui excitavi; quem non minus amo quam tu, paene dixi quam te. atque etiam spero toto anno imperi nostri terruncium sumptus in provincia nullum fore.

\ciceroSection{7} habes omnia. nunc publice litteras Romam mittere parabam. uberiores erunt quam si ex Amano misissem. at te Romae non fore! sed est totum in eo quid Kalendis Martiis futurum sit. vereor enim ne cum de provincia agetur, si Caesar resistet, nos retineamur. his tu si adesses, nihil timerem.

\ciceroSection{8} redeo ad urbana quae ego diu ignorans ex tuis iucundissimis litteris a. d. v Kal. Ianuarias denique cognovi. eas diligentissime Philogenes, libertus tuus, curavit perlonga et non satis tuta via perferendas. nam quas Laeni pueris scribis datas non acceperam. iucunda de Caesare et quae senatus decrevit et quae tu speras. quibus ille si cedit, salvi sumus. incendio Plaetoriano quod Seius ambustus es minus moleste fero. Lucceius de Q. Cassio cur tam vehemens fuerit et quid actum sit aveo scire.

\ciceroSection{9} ego cum Laodiceam venero, Quinto sororis tuae filio togam puram iubeor dare; cui moderabor diligentius. Deiotarus cuius auxiliis magnis usus sum ad me, ut scripsit, cum Ciceronibus Laodiceam venturus erat. tuas etiam Epiroticas exspecto litteras, ut habeam rationem non modo negoti verum etiam oti tui. Nicanor in officio est et a me liberaliter tractatur. quem , ut puto, Romam cum litteris publicis mittam, ut et diligentius perferantur et idem ad me certa de te et a te referat. Alexis quod mihi totiens salutem adscribit est gratum; sed cur non suis litteris idem facit quod meus ad te Alexis facit? Phemio quaeritur κέρασ . sed haec hactenus. cura ut valeas et ut sciam quando cogites Romam. etiam atque etiam vale.

\ciceroSection{10} tua tuosque Thermo et praesens Ephesi diligentissime commendaram et nunc per litteras ipsumque intellexi esse perstudiosum tui. tu velim, quod antea ad te scripsi, de domo Pammeni des operam ut quod tuo meoque beneficio puer habet cures ne qua ratione convellatur. utrique nostrum honestum existimo; tum mihi erit pergratum.
